\section{The margin is a sampling constant, and the density sets it}
\label{sec:sampling}

Section~\ref{sec:region} left the margin
$m_L=\lambda_{\min}(Q_{0,L})$ with an exact reading --- how well a function
supported in an interval can avoid the Riemann zeros in frequency --- and left its
collapse as the open mechanism question. This section identifies the object, proves
that it is positive, says what the classical theory of such objects reaches, and
then reports a measurement which changes what the mechanism question is.

\subsection{The object}

For $f\in C_c^\infty(I_L)$ the transform $\widehat F$ is entire of exponential
type $L/2$ and lies in $L^2(\R)$, that is $\widehat F\in PW_{L/2}$, the
Paley--Wiener space; and $\norm f^2=\frac1{2\pi}\norm{\widehat F}^2$. Under the
Riemann hypothesis the spectral form \eqref{eq:spectral-target} is a sum of
squares, so

\begin{equation}\label{eq:sampling}
 m_L=\lambda_{\min}(Q_{0,L})
 =\inf\Big\{\frac{2\pi\sum_\rho|\Phi(\gamma_\rho)|^2}{\norm\Phi^2}
 \ :\ \Phi\in PW_{L/2},\ \Phi\neq0\Big\} :
\end{equation}

\emph{the margin is the constant in the lower sampling inequality for the Riemann
zeros on $PW_{L/2}$.} Weil positivity says the zeros are a set of uniqueness for
$PW_{L/2}$; a positive $m_L$ says they are a set of sampling, with a bound. The
generating function of the sampling set is $\Xi(\tau)=\xi(\frac12+i\tau)$, a
Hermite--Biehler function under the hypothesis, which is why de Branges theory
keeps being named in this program \cite{SuzukiDeBranges}. Section~\ref{subsec:reach}
says how much that buys.

One half of the standard definition fails, and it is worth saying so before the
language is used further.

\begin{proposition}[No upper sampling constant]\label{prop:noupper}
Unconditionally, $\sup\{Q_{0,L}[f]:\norm f=1\}=+\infty$. Equivalently the zeros
are not relatively separated: $\#\{\rho:|\gamma_\rho-T|\leq\frac12\}\to\infty$.
\end{proposition}

\begin{proof}
By \eqref{eq:target} and \eqref{eq:multiplier}, $Q_{0,L}$ differs from
$\frac1{2\pi}\int|\widehat F|^2(b(\tau^2)+w_0)$ by the rank-two form $P_L$ and by
the finite prime comb, both bounded; and $b(\tau^2)+w_0=\log\frac{|\tau|}{2\pi}+
O(\tau^{-2})$ is unbounded. For the second statement, the number of zeros in a
unit window at height $T$ is $\frac1{2\pi}\log\frac T{2\pi}+o(\log T)$ by
Riemann--von Mangoldt.
\end{proof}

So $\{\gamma_\rho\}$ is \emph{not} a sampling set in the usual two-sided sense and
$m_L$ is \emph{not} a frame bound; it is the lower constant alone. This is the
same fact as the unboundedness of the group delay in
Proposition~\ref{prop:groupdelay}, and the same fact as the channel-tower rule of
\cite{Selection}. The lower inequality is the half that needs no separation, for
the trivial reason that $\sum_{\lambda\in\Lambda}|\Phi(\lambda)|^2$ increases when
$\Lambda$ does; that monotonicity is used again in \S\ref{subsec:density}.

\subsection{The infimum is positive}

\begin{proposition}[Strict positivity of the sampling constant]\label{prop:positive}
Assume the Riemann hypothesis. Then $m_L>0$ for every $L>0$, and the infimum in
\eqref{eq:sampling} is attained.
\end{proposition}

\begin{proof}
Write $Q_{0,L}=\mathcal A_L+\mathcal B_L$, where
$\mathcal A_L[f]=\frac1{2\pi}\int|\widehat F|^2(b(\tau^2)+w_0)$ and
$\mathcal B_L=P_L-(\text{prime comb})$ is bounded, of norm at most
$2\norm{\cosh(x/2)}^2_{L^2(I_L)}+2\sum_{n<e^L}\Lambda(n)n^{-1/2}$. Since
$b\geq0$ and $w_0=-5.37218\ldots$, $\mathcal A_L$ is bounded below, with form
domain $\mathcal D_{\log,L}$ of Definition~\ref{def:generator}.

\emph{The embedding $\mathcal D_{\log,L}\hookrightarrow L^2(I_L)$ is compact.}
Let $(f_k)$ be bounded in the form norm. Then $\widehat{F_k}$ is bounded in
$L^2(\R)$; its tails are uniformly small, $\int_{|\tau|>T}|\widehat{F_k}|^2\leq
C/\log T$; it is uniformly bounded, $|\widehat{F_k}|\leq\norm{f_k}_{L^1}\leq
\sqrt L\norm{f_k}$; and it is equicontinuous, since
$|\widehat{F_k}'|=|\widehat{xF_k}|\leq\frac{L^{3/2}}2\norm{f_k}$. Arzel\`a--Ascoli
and a diagonal argument give a subsequence converging locally uniformly, and the
uniform tail bound upgrades that to convergence in $L^2(\R)$.

Hence $\mathcal A_L$ has compact resolvent; adding the bounded $\mathcal B_L$
leaves the essential spectrum empty, so $Q_{0,L}$ has purely discrete spectrum
accumulating only at $+\infty$ and its infimum is attained at some $e_L$ with
$\norm{e_L}=1$. Under the hypothesis $Q_{0,L}\geq0$ on $C_c^\infty(I_L)$, hence on
the form closure. Suppose $m_L=0$. Along an approximating sequence
$\widehat{F_k}\to\widehat{E_L}$ locally uniformly, so for every $T$
\[
 \sum_{|\gamma_\rho|<T}|\widehat{E_L}(\gamma_\rho)|^2
 =\lim_k\sum_{|\gamma_\rho|<T}|\widehat{F_k}(\gamma_\rho)|^2
 \leq\liminf_kQ_{0,L}[f_k]=0 ,
\]
so $\widehat{E_L}$ vanishes at every zero. But $\widehat{E_L}\in PW_{L/2}$ is a
nonzero function of Cartwright class and type $L/2$, so the number of its zeros in
$[-T,T]$ is at most $\frac L\pi T(1+o(1))$, while the zeros of $\xi$ supply
$2N(T)=\frac T\pi\log\frac T{2\pi e}+\frac74+2S(T)$. The two are incompatible for
large $T$.
\end{proof}

Two remarks. The counting contradiction is the reason strict positivity holds and
it is a statement about \emph{densities}, not an estimate; attainment of the
minimum is Bombieri's \cite{BombieriWeil}, by a different argument, and the
continuity of $m_L$ in $L$ is Suzuki's \cite{Suzuki}. What the proof above adds is
the mechanism: the divergence of the group delay
(Proposition~\ref{prop:groupdelay}) is exactly what makes the form domain embed
compactly, so \textbf{the unboundedness of the Wigner--Smith delay and the
positivity of the sampling constant are one fact}. An all-pass filter with bounded
delay would have neither.

\subsection{The sampling budget, and where it crosses}

The local density of $\{\pm\gamma_\rho\}$ is $\frac1{2\pi}\log\frac\tau{2\pi}$ and
the Nyquist density for $PW_{L/2}$ is $\frac L{2\pi}$. They meet at
\begin{equation}\label{eq:crossover}
 \tau_c=2\pi e^L ,
\end{equation}
below which the zeros undersample and above which they oversample; the cumulative
counts cross at $T_*=e\,\tau_c=2\pi e^{L+1}$, which is where the proof of
Proposition~\ref{prop:positive} bites. The maximal shortfall is exact:
\begin{equation}\label{eq:deficit}
 D(L)=\frac{L\tau_c}\pi-2N(\tau_c)
 =\frac{\tau_c}\pi\Big(L-\log\frac{\tau_c}{2\pi e}\Big)-\frac74
 =2e^L-\frac74 ,
\end{equation}
the $\frac74$ being twice the Riemann--von Mangoldt constant and $S$ dropped. Since
$e^L=X$ is the prime cutoff in \eqref{eq:target}, \textbf{the sampling crossover is
the explicit-formula balance}: primes up to $X$ against zeros up to height $2\pi X$,
with a shortfall of $2X$ samples. This is the first place in the program where the
$e^L$ scale appears for a structural reason rather than as the range of a sum.

\subsection{What classical sampling theory reaches}\label{subsec:reach}

Very little, and it is useful to say exactly how little, because the language of
the subject is being used for an object outside its hypotheses.

Landau's necessary density condition is satisfied with room to spare, the lower
uniform density being infinite. Beurling's gap theorem is the only tool here that
produces a \emph{constant}; in the explicit form of \cite{OlevskiiUlanovskii} it
requires the maximal gap of the sampling set to be below $\pi/\sigma$ for type
$\sigma$. The maximal gap of $\{\pm\gamma_\rho\}$ is the one straddling the origin,
$2\gamma_1=28.2695$, and $\sigma=L/2$, so the hypothesis reads
\begin{equation}\label{eq:beurling}
 L<\pi/\gamma_1=0.2222606\ldots,
\end{equation}
with constant $(\cos(\gamma_1L/2))^{-1}$: $1.32$ at $L=0.1$ and $63$ at $L=0.22$.
\textbf{Above $L\approx0.22$ the classical machinery does not degrade; it stops
applying.} For comparison, Yoshida's positivity theorem reaches $L=\log2$ and
\cite{ZhuCompact} certifies $L=1.6$.

Ortega-Cerd\`a and Seip \cite{OrtegaCerdaSeip} characterize the sampling sequences
for $PW$ through de Branges theory --- a separated $\Lambda$ is sampling if and
only if it is the zero sequence of $EF+E^*F^*$ for suitable Hermite--Biehler
$E,F$ with $H(E)=PW$ --- but this is a dichotomy and produces no constant, as the
authors state. The one route in that literature which does carry constants, through
the $A_2$ condition on $|G|^2$ for the generating function $G$, is stated for
\emph{complete interpolating sequences}; the zeros are not one, since they
overshoot the critical density rather than matching it.

\begin{remark}[What Suzuki's de Branges space is, and is not]\label{rem:suzuki}
\cite{SuzukiDeBranges} shows that under RH the Hilbert space of the Weil
distribution is isomorphic to a de Branges space --- but with structure function
$E_\xi(z)=\xi(\frac12-iz)+\xi'(\frac12-iz)$, whose zero set is not the zeta zeros,
and the isomorphism is with the whole Weil space rather than with $PW_{L/2}$. Its
chain of subspaces is indexed by the support cut-off, which is structurally our
$L$, and its content is the conversion of Weil's inequality into an identity. It
contains no smallest eigenvalue, no gap and no rate. The two questions are
adjacent, not the same, and the de Branges side currently supplies nothing
quantitative. Earlier drafts of this program said otherwise.
\end{remark}

\subsection{Where the obstruction lives, and a construction that fails}
\label{subsec:profile}

Let $\lambda^{(W)}_L$ be the minimum of the Rayleigh quotient over trial functions
whose spectral content is confined below $W$. The fraction of the exponent it
recovers is, measured at six lengths:

\begin{center}
\begin{tabular}{@{}lcccccc@{}}
\toprule
$W/\tau_c$ & $0.50$ & $0.75$ & $1.00$ & $1.25$ & $1.50$ & $2.00$\\
\midrule
$L=1.6$ & $0.585$ & $0.733$ & $0.858$ & $0.933$ & $0.982$ & $0.995$\\
$L=2.0$ & $0.561$ & $0.725$ & $0.840$ & $0.910$ & $0.957$ & $0.995$\\
$L=2.2$ & $0.586$ & $0.744$ & $0.842$ & $0.921$ & $0.969$ & $0.997$\\
$L=2.5$ & $0.584$ & $0.739$ & $0.835$ & $0.915$ & $0.965$ & $0.996$\\
$L=2.8$ & $0.590$ & $0.744$ & $0.844$ & $0.924$ & $0.966$ & $0.997$\\
$L=3.0$ & $0.592$ & $0.741$ & $0.843$ & $0.915$ & $0.964$ & $0.997$\\
\bottomrule
\end{tabular}
\end{center}

\noindent
(entries $\log\lambda^{(W)}_L/\log m_L$). The columns are constant to about one
percent while $m_L$ moves through $10^{80}$: \textbf{the frequency profile of the
obstruction is universal in $W/\tau_c$}, so \eqref{eq:crossover} is the right
scale, and the obstruction is essentially complete at $2\tau_c$, between $\tau_c$
and $T_*$. There is a practical corollary, and it is the reason the numbers below
are reported with a basis-convergence rule rather than as single values: a trial
space must reach about $2.5\,\tau_c$, that is a dimension of about $5Le^L$, or the
computed value saturates and looks converged when it is not.

The construction the deficit \eqref{eq:deficit} suggests --- stay below $\tau_c$
and \emph{vanish} at the $2N(\tau_c)$ sub-critical zeros, leaving $D(L)$ free
dimensions --- does not work, and fails worse as $L$ grows:

\begin{center}
\begin{tabular}{@{}lrrrrr@{}}
\toprule
$L$ & $1.6$ & $2.0$ & $2.4$ & $2.8$ & $3.0$\\
\midrule
$\log_{10}m_L$ & $-16.75$ & $-29.17$ & $-48.04$ & $-76.75$ & $-96.25$\\
band $|\tau|<\tau_c$ only & $-14.30$ & $-24.46$ & $-40.41$ & $-64.84$ & $-81.19$\\
band and vanishing & $-14.30$ & $-24.46$ & $-36.49$ & $-45.62$ & $-50.64$\\
\midrule
fraction recovered & $0.85$ & $0.84$ & $0.76$ & $0.59$ & $0.53$\\
\bottomrule
\end{tabular}
\end{center}

The minimizer does not interpolate the undersampled zeros; it makes $\widehat F$
\emph{small at every zero}, including the oversampled ones above $\tau_c$, and
imposing exact vanishing is an expensive constraint. In the language of
time--frequency localization this construction is $1-\lambda_n$ for the prolate
operator at Shannon number $\mathcal N_L=2Le^L$ and index $\mathcal N_L-D(L)$,
hence at occupancy fraction $u_L=1-D(L)/\mathcal N_L$ bounded away from both
endpoints. That regime is neither the classical fixed-index one nor the
Landau--Widom plunge. \cite{ConnesConsaniProlate} already identifies prolate
eigenfunctions as the source of very small eigenvalues of the Weil form, without
a rate; the sharpest available statement about the rate is
$-\log(1-\lambda_n)\asymp(c-n)/\log\frac{2c}{c-n}$ with unequated constants, and
the corresponding limit at fixed occupancy is an open problem
\cite{Kulikov}.

\subsection{The collapse is set by the density, not by the arithmetic}
\label{subsec:density}

The natural control is a point set with the \emph{same counting function} and no
arithmetic. (Deleting the prime comb is not such a control: archimedean plus pole
is not a sum of squares over a point set, and is in fact indefinite for
$L\gtrsim0.85$.) We use the Gram points $g_n$, defined by $\theta(g_n)=n\pi$,
whose counting function is exactly the smooth part $\theta(T)/\pi+1$ with
$S\equiv0$; and jittered Gram points, the jitter applied in the unfolded variable
$\theta/\pi$ so that the counting function is preserved. Everything is computed in
the same basis, at the same truncation, and at the same precision.

\begin{center}
\begin{tabular}{@{}lrrr@{}}
\toprule
$L$ & $\log_{10}m_L$ (the zeros) & $\log_{10}$, Gram points & ratio of logarithms\\
\midrule
$0.6$ & $-2.11$ & $-1.00$ & $2.100$\\
$0.8$ & $-3.73$ & $-2.17$ & $1.715$\\
$1.0$ & $-6.01$ & $-4.06$ & $1.480$\\
$1.2$ & $-8.76$ & $-6.62$ & $1.323$\\
$1.4$ & $-12.35$ & $-9.97$ & $1.239$\\
$1.6$ & $-16.75$ & $-14.18$ & $1.181$\\
$1.8$ & $-22.35$ & $-19.49$ & $1.147$\\
$2.0$ & $-29.17$ & $-26.15$ & $1.116$\\
$2.2$ & $-37.62$ & $-34.48$ & $1.091$\\
$2.4$ & $-48.04$ & $-44.77$ & $1.073$\\
$2.6$ & $-60.97$ & $-57.48$ & $1.061$\\
$2.8$ & $-76.75$ & $-73.17$ & $1.049$\\
\bottomrule
\end{tabular}
\end{center}

At $L=2$ the four arithmetic-free sets --- Gram, and jitter of $0.15$, $0.40$ and
$1.00$ of a mean spacing --- give $\log_{10}$ of $-26.15$, $-25.99$, $-26.53$ and
$-25.44$ against $-29.17$ for the zeros; at $L=1.6$, $-14.18$, $-14.02$, $-14.59$
and $-13.50$ against $-16.75$. Three readings.

\begin{enumerate}
\item \textbf{Every arithmetic-free set of the same density collapses
superexponentially, at essentially the same rate.} The ratio of logarithms falls
monotonically from $2.100$ at $L=0.6$ to $1.049$ at $L=2.8$.
\item \textbf{The arithmetic is worth about three orders of magnitude and grows
like $\log L$}, against a total growing like $e^L$. It is a systematic bias and not
noise: jitter of up to a full mean spacing moves $\log_{10}m_L$ by about one in
either direction, and the zeros sit below the whole arithmetic-free family.
\item \textbf{So the margin carries almost no arithmetic information.} Its collapse
is what the density $\frac1{2\pi}\log\frac\tau{2\pi}$ does to the lower sampling
constant of a Paley--Wiener space whose type is far below it.
\end{enumerate}

\begin{remark}[Scope of the measurement]\label{rem:densityscope}
The values for the control sets are trial-space Rayleigh quotients with the sum
truncated, hence upper bounds on their own infima and slightly depressed by
truncation, so the comparison is conservative in the direction that matters: the
control values could only be higher, which would widen the gap without changing
its relative decline. The zeros row is computed from the explicit formula without
truncation. One jitter draw per strength; the spread across strengths is the error
bar quoted, not a statistical estimate. The range is $L\in[0.6,3.0]$, which is
$96$ orders of magnitude in $m_L$ but only a factor of five in $L$, and the ratio
of logarithms has not been shown to converge.
\end{remark}

\subsection{What this changes}

Section~\ref{sec:region} ended by naming the collapse of the zero-placement
precision as the mechanism question. Proposition~\ref{prop:kappa} stands and
\S\ref{subsec:density} does not contradict it; what it does is remove the reason to
ask. A Gram-point analogue of $\xi$, with no primes anywhere in it, has a margin
that collapses at the same rate to within a ratio of logarithms tending to one.

\begin{remark}[Consequence for the program]\label{rem:mechanism}
Two things follow, and they point in opposite directions.
Reproducing the superexponential collapse is \emph{weak evidence} for a candidate
source, because the density alone produces it; and failing to reproduce it is weak
evidence against one. But a candidate is correspondingly \emph{not} required to
produce an arithmetic rate. It has to produce a density --- a group delay growing
like $\log\tau$, which is part of Condition~\ref{cond:bilinear} --- and that is a far cheaper
demand than a superexponential margin law. The quantity that does carry arithmetic
is the \emph{gap} between the zeros and matched-density sets, about $10^3$ at
$L=2$; nothing in this program or in the literature has studied it.
\end{remark}

\begin{remark}[On decay laws, and a convention]\label{rem:laws}
No one-parameter law in $D(L)$, in $\mathcal N_L$, in $\mathcal N_L/\log\mathcal N_L$,
in $N(\tau_c)/\log N(\tau_c)$ or in $D/\log(2\mathcal N_L/D)$ fits
$L\in[1.4,3.0]$ better than $22\%$ in the exponent. Over any factor-$1.25$ window
in $L$ all of them fit to a few percent, which is why fitting on three or four
points determines nothing. Two published fits illustrate this: the relation
recorded in \cite{CriticalPath}, refitted on its own three horizons, predicts
$L=3.0$ to within $4.1\%$ and $L=4.0$ to within $6.5\%$; the Landau--Widom form of
\cite{ZhuCompact}, fitted on $L\in[2.8,4.0]$, is $46\%$ high at $L=1.6$.

A convention warning, because it has already cost this program one wrong
conclusion. \cite{ZhuCompact} supports $f$ in $[-L,L]$ and sums primes over
$\log n<2L$; this manuscript supports $f$ in $(-L/2,L/2)$ and sums over $n<e^L$.
\textbf{The $L$ of \cite{ZhuCompact} is half the $L$ used here.} With the
conversion applied, its certified enclosure $8.9\times10^{-18}\leq m\leq
2.27\times10^{-17}$ belongs at our $L=1.6$, where the value computed here is
$1.73\times10^{-17}$, and its $3.19\times10^{-283}$ at our $L=4$.
\end{remark}

\begin{remark}[Status of the numbers in this section]\label{rem:samplingstatus}
Propositions~\ref{prop:noupper} and \ref{prop:positive} are written proofs, the
second conditional on RH and quoting the Cartwright density theorem.
Equations~\eqref{eq:crossover}--\eqref{eq:beurling} are written computations,
verified in Appendix~\ref{app:checks}. Every table is a labelled numerical
computation, produced by unregistered \texttt{mpmath} programmes recorded in
\cite{FrameBoundNote}; each entry is a trial-space Rayleigh quotient and hence a
one-sided bound, and the direction is stated where a comparison is made. The
assembly those programmes use agrees with an independent quadrature assembler of
this program to a relative $6\times10^{-25}$, and reproduces the three margins of
\cite{CriticalPath} to $2.2\%$, $3.2\%$ and $0.2\%$ --- which, incidentally, puts
Proposition~\ref{prop:kappa} on firmer ground than the three data points it was
stated with. Section~\ref{subsec:reach} is a reading of the cited literature.
\end{remark}
